\chapter*{Appendix N: Rewriting Systems and Normal Forms}
\addcontentsline{toc}{chapter}{Appendix N: Rewriting Systems and Normal Forms}

\section{Abstract Rewriting Systems}

\subsection{Definition}

An abstract rewriting system (ARS) is a pair $(A,\to)$ where:
\begin{itemize}
\item $A$ is a set of objects,
\item $\to \subseteq A \times A$ is a binary relation called the rewrite relation.
\end{itemize}

\subsection{Rewrite Sequences}

A finite rewrite sequence is a chain
\[
a_0 \to a_1 \to \cdots \to a_n.
\]
A transitive closure is denoted $\to^+$, and the reflexive–transitive closure by $\to^\ast$.

\section{Termination and Confluence}

\subsection{Termination}

The system is terminating if there is no infinite sequence
\[
a_0 \to a_1 \to a_2 \to \cdots.
\]

\subsection{Confluence}

The system is confluent if whenever $a \to^\ast b$ and $a \to^\ast c$, there exists $d$ such that
\[
b \to^\ast d
\quad\text{and}\quad
c \to^\ast d.
\]

\subsection{Local Confluence}

Local confluence holds if for every peak
\[
a \leftarrow b \to c
\]
there exists $d$ with $b \to^\ast d$ and $c \to^\ast d$.

\subsection{Newman's Lemma}

If an ARS is terminating and locally confluent, then it is confluent.

\section{Normal Forms}

\subsection{Definition}

An element $n\in A$ is a normal form if there is no $a$ with $n\to a$.

\subsection{Unique Normal Forms}

A system has unique normal forms if:
\begin{itemize}
\item every $a\in A$ reduces to some normal form $n$,
\item if $a\to^\ast n_1$ and $a\to^\ast n_2$ with $n_1,n_2$ normal forms, then $n_1 = n_2$.
\end{itemize}

Confluence implies unique normal forms.

\section{Critical Pairs}

\subsection{Definition}

Let $R$ be a set of rewrite rules.  
A critical pair is a pair of terms $(s,t)$ arising from overlapping applications of rules in $R$.

\subsection{Joinability}

A critical pair $(s,t)$ is joinable if there exists $u$ such that
\[
s \to^\ast u
\quad\text{and}\quad
t \to^\ast u.
\]

\subsection{Critical Pair Theorem}

If the system is terminating and all critical pairs are joinable, then it is confluent.

\section{Orthogonality}

\subsection{Left-Linearity}

A rewrite rule is left-linear if no variable appears more than once in the left-hand side.

\subsection{Non-Overlapping}

A rule set is non-overlapping if no left-hand sides overlap except by variable substitution.

\subsection{Orthogonality Theorem}

If a TRS is orthogonal (left-linear and non-overlapping), then it is confluent.

\section{Contextual Rewriting}

\subsection{Contexts}

A context is a term with a distinguished hole $[\,]$.  
If $C$ is a context and $t$ a term, the hole-filling operation is $C[t]$.

\subsection{Contextual Closure}

The rewrite relation is closed under context:
\[
s\to t \ \Rightarrow\ C[s] \to C[t].
\]

\section{Rewriting Modulo Equations}

\subsection{Definition}

Let $\equiv$ be a congruence on $A$.  
Rewriting modulo $\equiv$ means:
\[
a \to b \quad\text{iff there exist } a',b' \text{ such that }
a \equiv a',\ a'\to b',\ b'\equiv b.
\]

\subsection{Properties}

If the base system is terminating and $\equiv$ is decidable and confluent, then rewriting modulo $\equiv$ preserves termination and confluence.

\section{Completion Procedures}

\subsection{Knuth--Bendix Completion}

Given a terminating but non-confluent system, completion attempts to generate new rewrite rules to ensure joinability of all critical pairs.

\subsection{Outcome}

The procedure yields:
\begin{itemize}
\item a confluent system extending the original rules, or
\item a failure if no finite completion exists.
\end{itemize}

\section{Rewriting on Graphs}

\subsection{Hypergraph Rewriting}

A rule consists of:
\begin{itemize}
\item a pattern hypergraph $L$,
\item a replacement hypergraph $R$,
\item an injective match $m:L\hookrightarrow H$ into a host hypergraph $H$.
\end{itemize}

The rewritten graph is obtained by removing the image of $L$ and gluing in $R$ along the boundary.

\subsection{Pushout-Based Rewriting}

A rewriting step is defined by the double-pushout construction:
\[
L \leftarrow K \rightarrow R,
\]
where $K$ is the interface graph.

\section{Normal Forms in Graph Rewriting}

\subsection{Graph Normality}

A hypergraph $H$ is a normal form if no rule applies to $H$.

\subsection{Uniqueness}

If all matches are monomorphic and the rule set satisfies the gluing condition, unique normal forms are obtained.

\section{Rewriting Metrics}

\subsection{Rewrite Distance}

The rewrite distance between $a$ and $b$ is:
\[
d(a,b) =
\min\{ n \mid a \to^n b' \ \text{and}\ b' = b \}.
\]

\subsection{Energy-Decreasing Rewriting}

If a functional $\mathcal{E}:A\to\mathbb{R}_{\ge 0}$ satisfies  
$\mathcal{E}(b) < \mathcal{E}(a)$ for each $a\to b$,  
then termination follows.

\section{Completion of Rewriting Systems}

\subsection{Canonical Systems}

A rewriting system is canonical if it is:
\begin{itemize}
\item terminating,
\item confluent,
\item and has a unique normal form for every element.
\end{itemize}

\subsection{Canonical Form Function}

For such a system, the function
\[
\mathrm{nf} : A \to A
\]
mapping each element to its unique normal form is well-defined, idempotent, and satisfies:
\[
a \to^\ast \mathrm{nf}(a).
\]


